\documentclass[12pt]{article}
\usepackage[utf8]{inputenc}
\usepackage{geometry}
\geometry{letterpaper, margin=0.25in}
\usepackage{graphicx} 
\usepackage{multicol}
\usepackage{parskip}
\usepackage{booktabs}
\usepackage{array} 
\usepackage{paralist} 
\usepackage{verbatim}
\usepackage{sectsty}
\usepackage[shortlabels]{enumitem}

\usepackage{amsmath}
\usepackage{amssymb}
\usepackage{mathtools}
\usepackage{empheq}
\usepackage{xcolor}
\usepackage{bbm}
\usepackage{tikz}
\usepackage{pgfplots}
\usepackage{tikz-cd}
\usepackage{circuitikz}
\usepackage{multirow}
\pgfplotsset{compat=1.18}
\usetikzlibrary{intersections, decorations.markings}
\tikzset{
    marking along/.style n args={2}{
        decoration={
                markings, 
                mark=at position #1 with {\arrow{#2}}
        },
        postaction={decorate}
        },
    marking along/.default={0.5}{>}
    wavy/.style={
        decorate,decoration={coil,aspect=0}
     },
    two marks/.style n args={1}{
        decoration={
            markings,
            mark=at position 0.25 with {\arrow{#1}},
            mark=at position 0.75 with {\arrow{#1}}
         },
         postaction={decorate}
    }
    two marks/.default={>}
}

\colorlet{mygreen}{green!50!teal}
\colorlet{darkblue}{blue!60!black}

\newcommand{\ans}[1]{\boxed{\text{#1}}}
\newcommand{\vecs}[1]{\langle #1\rangle}
\renewcommand{\hat}[1]{\widehat{#1}}

\renewcommand{\P}{\mathbb{P}}
\newcommand{\R}{\mathbb{R}}
\newcommand{\E}{\mathbb{E}}
\newcommand{\Z}{\mathbb{Z}}
\newcommand{\N}{\mathbb{N}}
\newcommand{\Q}{\mathbb{Q}}
\newcommand{\C}{\mathbb{C}}

\newcommand{\ind}{\mathbbm{1}}
\newcommand{\qed}{\quad \blacksquare}

\newcommand{\brak}[1]{\left\langle #1 \right\rangle}
\newcommand{\bra}[1]{\left\langle #1 \right\vert}
\newcommand{\ket}[1]{\left\vert #1 \right\rangle}

\newcommand{\abs}[1]{\left\vert #1 \right\vert}
\newcommand{\mfX}{\mathfrak{X}}
\newcommand{\ep}{\varepsilon}

\newcommand{\Ec}{\mathcal{E}}
\newcommand{\Nc}{\mathcal{N}}
\newcommand{\A}{\mathcal{A}}
\newcommand{\F}{\mathcal{F}}
\newcommand{\Cc}{\mathcal{C}}
\newcommand{\B}{\mathcal{B}}
\newcommand{\M}{\mathcal{M}}
\newcommand{\X}{\chi}
\renewcommand{\L}{\mathcal{L}}

\newcommand{\sub}{\subseteq}
\newcommand{\st}{\text{ s.t. }}
\newcommand{\card}{\text{card }}
\renewcommand{\div}{\vspace*{10pt}\hrule\vspace*{10pt}}
\newcommand{\surj}{\twoheadrightarrow}
\newcommand{\inj}{\hookrightarrow}
\newcommand{\biject}{\hookrightarrow \hspace{-8pt} \rightarrow}
\renewcommand{\bar}[1]{\overline{#1}}
\newcommand{\overcirc}[1]{\overset{\circ}{#1}}
\newcommand{\diam}{\text{diam }}
\newcommand{\iid}{\overset{	ext{iid}}{\sim}}

\renewcommand{\Re}{\text{Re}\,}
\renewcommand{\Im}{\text{Im}\,}
\newcommand{\Var}{\text{Var}\,}
\newcommand{\Cov}{\text{Cov}\,}

\DeclareMathOperator*{\argmax}{\arg\max}
\DeclareMathOperator*{\argmin}{\arg\min}

\newcommand{\sign}{\text{sign}\,}

\newcommand*{\tbf}[1]{\ifmmode\mathbf{#1}\else\textbf{#1}\fi}

\usepackage{tcolorbox}
\tcbuselibrary{breakable, skins}
\tcbset{enhanced}
\newenvironment*{tbox}[2][gray]{
    \begin{tcolorbox}[
        parbox=false,
        colback=#1!5!white,
        colframe=#1!75!black,
        breakable,
        title={#2}
    ]}
    {\end{tcolorbox}}

\newenvironment*{exercise}[1][red]{
    \begin{tcolorbox}[
        parbox=false,
        colback=#1!5!white,
        colframe=#1!75!black,
        breakable
    ]}
    {\end{tcolorbox}}

\newenvironment*{proof}[1][blue]{
\begin{tcolorbox}[
    parbox=false,
    colback=#1!5!white,
    colframe=#1!75!black,
    breakable
]}
{\end{tcolorbox}}
\newenvironment*{proposition}[1][gray]{
    \begin{tcolorbox}[
        parbox=false,
        colback=#1!5!white,
        colframe=#1!75!black,
        breakable
    ]}
    {\end{tcolorbox}}


\begin{document}
\begin{multicols}{2}
    \begin{flushleft}
        \parbox{0.5\textwidth}{\Large ENGN 1580 - Homework 3 }
    \end{flushleft}
    \begin{flushright}
        \tbf{Milan Capoor}\\
        Spring 2026
    \end{flushright}
\end{multicols}
\div

\tbf{1. Convexity.} Is the function $f (x_1, x_2) = (x_1 + x_2)^2$ convex?

\color{darkblue}
    \[H = \begin{pmatrix}
        \frac{\partial^2 f}{\partial x_1^2} & \frac{\partial^2 f}{\partial x_1 \partial x_2}\\ 
        \frac{\partial^2 f}{\partial x_2 \partial x_1} & \frac{\partial^2 f}{\partial x_2^2}
    \end{pmatrix} = \begin{pmatrix}
        2 & 2\\ 
        2 & 2    
    \end{pmatrix} \implies \lambda \in \{0, 2\} \implies f \text{ convex}\]

\color{black}

\tbf{2. Strict Convexity: } Is the function $f (x_1, x_2) = (x_1 + x_2)^2 $strictly convex?

\color{darkblue}
    $\lambda = 0$ so not strictly convex. 
\color{black}

\tbf{3. Optimization:} A convex region is defined similarly to a convex function:
\begin{itemize}
\item pick two points
\item draw a line between them
\item if all the points on the line are contained in the region, the region is convex
\item if not, then it’s not convex
\end{itemize}

\begin{enumerate}[a.]
    \item Please show that when minimizing a strictly convex function $f (\cdot)$ over its entire domain
    that the minimizing point is unique.

    \color{darkblue}
        Suppose $f(x) = f(y) = \min f(\cdot)$ but $x \neq y$. Then 
        \[f(\lambda x + (1 - \lambda) y) < \lambda f(x) + (1 - \lambda)f(y) = f(y)\]
        but this contradicts minimality of $y$. 
    \color{black}


    \item Please show that the minimum of a strictly convex $f (x)$ on a non-convex region might
    not be unique.
    \item If a continuous function $f (x)$ is strictly convex over $x \in \R^n$ , then does the absolute minimum occur at the point where all the N first partials are zero?
    
    \color{darkblue}
        Yes. 
    \color{black}


    \item If a function is NOT convex, can we be assured the minimum occurs when the first
    partials are zero?

    \color{darkblue}
        No. Consider $f(x) = x^3(x - 2)$. Not convex and global minimum is at $x = 2$ but $f_x(0) = 4x^2 - 6x^2 = 0$. 
    \color{black}

\end{enumerate}

\tbf{4. Quantization and Convexity:} You are designing a 2-level quantizer with parameters $x_0, q_0, q_1$.
\begin{enumerate}[a.]
    \item Please derive the Hessian matrix for the three design parameters.
    
    \color{darkblue}
        \begin{align*}
            Q(x) &= \begin{cases}
                q_0 & x < x_0\\ 
                q_1 & x \geq x_0
            \end{cases} = q_0 \ind_{x \in [-\infty, x_0]}(x) + q_1 \ind_{x \in [0, \infty]}(x) \\ 
            H_Q &= \begin{pmatrix}
                1
            \end{pmatrix}
        \end{align*}
        
    \color{black}

    \item Is the Hessian matrix positive definite? Positive semi-definite?
\end{enumerate}
\end{document}